\chapter{Protocolo experimental}
\label{cap:protocolo}

% GUION. Fuentes: src/particion.py, doc/plan_eda.md, doc/memoria/materiales.md seccion 5.
% ATENCION: el esqueleto anterior hablaba de validacion en hardware Eagle/Osprey.
% Eagle esta RETIRADA y el plan abierto da 10 min/mes: esa validacion no forma parte
% del alcance vigente. Si se menciona, que sea como trabajo futuro.

\section{La partición}
\label{sec:particion}
% Fuente: src/particion.py (modulo compartido por los tres modelos)
\subsection{Por qué la validación es temporal y no aleatoria}
El test es OOD en el tiempo: sus días de calibración son posteriores y disjuntos. Una
validación aleatoria compartiría días con el entrenamiento y daría una métrica
\textbf{optimista que no anticipa} lo que pasará en test.

Reservando los últimos días se obtiene una validación que \textbf{imita la relación real}
entre entrenamiento y test. Es más pequeña y más difícil, y esa es justamente la gracia.

\subsection{Una sola definición para los tres modelos}
La partición vive en un \textbf{único módulo} que importan tanto la vía tabular como el
cargador de grafos. \textbf{Es lo que convierte «protocolo idéntico» en algo demostrable} en
vez de en una afirmación.

Comprobación explícita: ningún día se comparte entre entrenamiento, validación y test.

\section{El preprocesado}
\label{sec:preprocesado}
% Fuente: las 17 directrices de notebooks/eda/04_directrices.ipynb seccion 4
Cada decisión sale de una medida del análisis exploratorio, no del criterio:
\begin{itemize}
    \item \textbf{normalización} obligatoria para Ridge (las variables abarcan diez órdenes de
          magnitud), innecesaria para Random Forest;
    \item \textbf{logaritmos}, porque la fidelidad cae de forma multiplicativa con el número
          de puertas;
    \item \textbf{poda} de las dimensiones constantes o duplicadas;
    \item \textbf{envoltura de ángulos}, porque una rotación y la misma más una vuelta
          completa son la misma operación física;
    \item \textbf{exclusión del día}, porque cae fuera del rango visto el 100\,\% del tiempo.
\end{itemize}

\textbf{Regla que no se negocia:} todo se ajusta \textbf{solo con el conjunto de
entrenamiento}. Normalizar antes de particionar introduciría información del test.

\section{Métricas}
\label{sec:metricas}
MAE, RMSE, $R^2$ y mejora relativa frente a no mitigar.

\subsection{Cómo hay que reportarlas}
% Fuente: las 17 directrices. Varias salen de trampas detectadas.
\begin{itemize}
    \item \textbf{desglosadas por los tres ejes OOD} y por observable — un promedio global
          quedaría dominado por un solo tipo de circuito;
    \item \textbf{mediana y percentiles además de la media}, porque las distribuciones tienen
          colas muy pesadas;
    \item \textbf{también restringidas a las muestras con señal}, porque una fracción grande
          del conjunto no tiene nada que predecir e infla artificialmente cualquier promedio;
    \item en el eje de tamaño, \textbf{acompañadas de una métrica relativa}, porque el MAE baja
          solo porque el objetivo se encoge.
\end{itemize}

\subsection{El listón}
\textbf{No mitigar} es el baseline que da sentido al trabajo. Un modelo que no lo bata no
sirve, por bueno que sea su $R^2$.

\section{Reproducibilidad y trazabilidad}
\label{sec:mlops}
Semillas globales, versionado del conjunto de datos, registro de experimentos con MLflow, y el
manifiesto que acompaña al conjunto con versiones y configuración.

\section{Infraestructura de cómputo}
\label{sec:infra}
Generación en máquinas virtuales de Google Cloud —necesaria, no cómoda: las muestras de 15
qubits piden 17,2\,GB por proceso—. Entrenamiento del \textit{Graph Transformer} en GPU.
